\subsection{Resultados da quadratura de Chebyshev--Lobatto}
\label{subsec:cheb_lobatto_results}

A Tabela~\ref{tab:cheb_lobatto_results} reproduz os valores numéricos obtidos ao
executar o script \texttt{exercise01\_cheb\_lobatto.jl} com a quadratura
Chebyshev--Lobatto após a correção dos pesos. Cada linha apresenta a ordem $n$,
o valor computado, a referência analítica e os erros absoluto e relativo.

\begin{table}[h]
    \centering
    \caption{Resultados da quadratura de Chebyshev--Lobatto}
    \label{tab:cheb_lobatto_results}
    \begin{tabular}{lrrrrr}
        \toprule
        integrando & $n$ & numérico & exato & erro abs. & erro rel. \\
        \midrule
        $f_1(x)=1$ & 10 & 2{,}00000000 & 2{,}00000000 & 2{,}22\times10^{-16} & 1{,}11\times10^{-16} \\
        $f_1(x)=1$ & 20 & 2{,}00000000 & 2{,}00000000 & 4{,}44\times10^{-16} & 2{,}22\times10^{-16} \\
        $f_1(x)=1$ & 40 & 2{,}00000000 & 2{,}00000000 & 0{,}00\times10^{0} & 0{,}00\times10^{0} \\
        $f_1(x)=1$ & 80 & 2{,}00000000 & 2{,}00000000 & 2{,}22\times10^{-15} & 1{,}11\times10^{-15} \\
        $f_1(x)=1$ & 160 & 2{,}00000000 & 2{,}00000000 & 3{,}91\times10^{-14} & 1{,}95\times10^{-14} \\
        $f_2(x)=x^2$ & 10 & 0{,}66666667 & 0{,}66666667 & 3{,}33\times10^{-16} & 4{,}99\times10^{-16} \\
        $f_2(x)=x^2$ & 20 & 0{,}66666667 & 0{,}66666667 & 3{,}33\times10^{-16} & 4{,}99\times10^{-16} \\
        $f_2(x)=x^2$ & 40 & 0{,}66666667 & 0{,}66666667 & 4{,}44\times10^{-16} & 6{,}66\times10^{-16} \\
        $f_2(x)=x^2$ & 80 & 0{,}66666667 & 0{,}66666667 & 2{,}11\times10^{-15} & 3{,}16\times10^{-15} \\
        $f_2(x)=x^2$ & 160 & 0{,}66666667 & 0{,}66666667 & 4{,}52\times10^{-14} & 6{,}78\times10^{-14} \\
        $f_3(x)=(2x^2-1)^2$ & 10 & 0{,}93333333 & 0{,}93333333 & 1{,}78\times10^{-15} & 1{,}90\times10^{-15} \\
        $f_3(x)=(2x^2-1)^2$ & 20 & 0{,}93333333 & 0{,}93333333 & 2{,}22\times10^{-16} & 2{,}38\times10^{-16} \\
        $f_3(x)=(2x^2-1)^2$ & 40 & 0{,}93333333 & 0{,}93333333 & 9{,}99\times10^{-16} & 1{,}07\times10^{-15} \\
        $f_3(x)=(2x^2-1)^2$ & 80 & 0{,}93333333 & 0{,}93333333 & 6{,}66\times10^{-16} & 7{,}14\times10^{-16} \\
        $f_3(x)=(2x^2-1)^2$ & 160 & 0{,}93333333 & 0{,}93333333 & 1{,}63\times10^{-13} & 1{,}75\times10^{-13} \\
        $f_4(x)=\sqrt{1-x^{2}}$ & 10 & 1{,}56989859 & 1{,}57079633 & 8{,}98\times10^{-4} & 5{,}72\times10^{-4} \\
        $f_4(x)=\sqrt{1-x^{2}}$ & 20 & 1{,}57068679 & 1{,}57079633 & 1{,}10\times10^{-4} & 6{,}97\times10^{-5} \\
        $f_4(x)=\sqrt{1-x^{2}}$ & 40 & 1{,}57076029 & 1{,}57079633 & 3{,}60\times10^{-5} & 2{,}29\times10^{-5} \\
        $f_4(x)=\sqrt{1-x^{2}}$ & 80 & 1{,}57328578 & 1{,}57079633 & 2{,}49\times10^{-3} & 1{,}58\times10^{-3} \\
        $f_4(x)=\sqrt{1-x^{2}}$ & 160 & 1{,}57087560 & 1{,}57079633 & 7{,}93\times10^{-5} & 5{,}05\times10^{-5} \\
        \bottomrule
    \end{tabular}
\end{table}

\paragraph{Comparação com as integrais analíticas.}
Os três polinômios ($f_1$, $f_2$ e $f_3$) agora são integrados com erro de
arredondamento, confirmando que a quadratura reproduz exatamente todos os
monômios até ordem $n$. Para $f_4(x)=\sqrt{1-x^2}$, que não é polinomial,
observa-se convergência rápida: o erro absoluto cai para a faixa
$10^{-5}$--$10^{-4}$ quando $n$ cresce, com oscilações inerentes aos coeficientes
obtidos pelo sistema linear em ordens elevadas.

\paragraph{Explicação.}
As primeiras execuções utilizaram os pesos ``$w_0=w_n=\pi/(2n)$'' e
``$w_k=\pi/n$'' derivados da quadratura com peso $1/\sqrt{1-x^2}$; como a soma
destes pesos é $\pi$, a regra aproximava $\int f(x)/\sqrt{1-x^2}\,\mathrm{d}x$,
superestimando integrais com respeito a $\mathrm{d}x$. Para obter pesos válidos
no sentido usual, resolvemos um sistema linear tipo Vandermonde: construímos a
matriz $A_{ik}=x_k^i$ com os nós de Chebyshev--Lobatto e impomos as identidades
dos momentos $\int_{-1}^{1} x^i \mathrm{d}x = 2/(i+1)$ (quando $i$ é par) ou
zero (quando $i$ é ímpar). O vetor de pesos $w = A^{-1} b$ garante que qualquer
polinômio de grau $\le n$ seja integrado exatamente. Essa abordagem elimina o
fator $\pi$ indevido e produz uma quadratura consistente com a medida de
Lebesgue em $[-1,1]$.

\subsection{Implementação e fluxo do exercício}

As funções foram adicionadas ao módulo
\texttt{SpectralTools.ChebyshevLobattoQuadrature}. O código a seguir resume a
construção dos nós, dos pesos via sistema linear e da rotina de integração:

\begin{verbatim}
@inline function cheb_lobatto_nodes(n::Integer)
    N = _checked_order(n)
    nodes = Vector{Float64}(undef, N + 1)
    @inbounds for k in 0:N
        nodes[k + 1] = cospi(k / N)
    end
    return nodes
end

@inline function cheb_lobatto_weights(n::Integer)
    nodes = cheb_lobatto_nodes(n)
    A = Matrix{Float64}(undef, n + 1, n + 1)
    @inbounds for i in 0:n
        A[i + 1, :] = nodes .^ i
    end
    b = [iseven(i) ? 2.0 / (i + 1) : 0.0 for i in 0:n]
    return A \ b
end

@inline function cheb_lobatto_quadrature(f::Function, n::Integer)
    nodes = cheb_lobatto_nodes(n)
    weights = cheb_lobatto_weights(n)
    return sum(weights .* f.(nodes))
end
\end{verbatim}

O script \texttt{exercise01\_cheb\_lobatto.jl} instancia os quatro integrandos e
sua lista de integrais exatas, percorre os valores de $n$ e grava os resultados:

\begin{verbatim}
const QUAD_ORDERS = [10, 20, 40, 80, 160]
const EXERCISE_FUNCTIONS = [
    ("f₁(x) = 1", x -> 1.0, 2.0),
    ("f₂(x) = x^2", x -> x^2, 2 / 3),
    ("f₃(x) = (2x^2 - 1)^2", x -> (2x^2 - 1)^2, 14 / 15),
    ("f₄(x) = sqrt(1 - x^2)", x -> sqrt(1 - x^2), pi / 2),
]

for (label, f, exact) in EXERCISE_FUNCTIONS
    for n in QUAD_ORDERS
        approx = cheb_lobatto_quadrature(f, n)
        abs_err = abs(approx - exact)
        rel_err = abs_err / abs(exact)
        # impressão e armazenamento em DataFrame
    end
end
\end{verbatim}

Cada linha calculada é enviada a um \texttt{DataFrame}; ao final, o script imprime
as tabelas para inspeção e salva o arquivo
\texttt{results\_exercise01\_cheb\_lobatto.csv}, documento de referência para a
subseção de análise de resultados.

\subsection{Instalação do pacote e preparação do ambiente}

\paragraph{1. Instalação do Julia.}
\begin{enumerate}
    \item Baixe e instale o \texttt{juliaup}, disponível em
    \url{https://github.com/JuliaLang/juliaup}.
    \item Em sistemas Unix, a instalação pode ser feita com:
    \texttt{curl -fsSL https://install.julialang.org | sh}.
    \item Após a instalação, confirme executando \texttt{juliaup status}
    e \texttt{julia +release} (ou a versão específica desejada).
\end{enumerate}

\paragraph{2. Clonar o repositório e ativar o pacote.}
\begin{enumerate}
    \item Clone o repositório oficial:
    \texttt{git clone https://github.com/stealth-lndrs/SpectralTools.git}.
    \item Entre no diretório: \texttt{cd SpectralTools}.
    \item Abra o REPL do Julia: \texttt{julia}.
    \item Ative o projeto: no modo \texttt{pkg} (tecla \texttt{]}), execute
    \texttt{activate .} e depois \texttt{instantiate}.
    \item Para desenvolver o pacote em outro projeto (como o exercício),
    utilize \texttt{pkg> dev /caminho/para/SpectralTools}.
\end{enumerate}

\paragraph{3. Garantir dependências e testes.}
\begin{enumerate}
    \item Ainda no modo \texttt{pkg}, execute \texttt{test} para rodar a suíte
    de testes do pacote e validar a instalação.
    \item Para atualizar dependências, use \texttt{pkg> update}, seguido de
    \texttt{pkg> test}.
    \item Recomenda-se utilizar ambientes dedicados (\texttt{--project=.}) ao
    rodar scripts que dependem do \texttt{SpectralTools}.
\end{enumerate}

Esses passos garantem que qualquer usuário, mesmo sem o Julia previamente
instalado, consiga preparar o ambiente, instalar o pacote do GitHub e
reproduzir os resultados apresentados.
